\documentclass[border=10pt]{standalone}
\usepackage{tidal-tikz-styles}

\begin{document}
\begin{tikzpicture}[
    node distance=0.45cm and 0.8cm,
    every node/.append style={text opacity=1},
]

% ===================================================================
%  TITLE
% ===================================================================
\node[font=\normalsize\bfseries, text=black] (title)
    {Lagrangian Linearisation Pipeline};
\node[below=0.02cm of title, font=\tiny\sffamily, text=GrayLine]
    {Full nonlinear $\mathcal{L}$ $\to$ linearised EOM via xPert};

% ===================================================================
%  ENTRY
% ===================================================================

\node[stage=Gray, below=0.5cm of title] (entry)
    {\textbf{Input}: full Lagrangian $\mathcal{L}(g, \phi, A, \ldots)$\\[-1pt]
     \scriptsize metric + matter fields, constants, background config};

% ===================================================================
%  STEP 1: Second-order perturbation
% ===================================================================

\node[stage=Purple, below=0.5cm of entry, text width=4.6cm] (step1)
    {\textbf{Step 1}: Quadratic Lagrangian\\[-1pt]
     \scriptsize $\mathcal{L}^{(2)} =
     \frac{1}{2}\frac{\delta^2(\sqrt{|g|}\,\mathcal{L})}
     {\sqrt{|g_0|}}$};

\node[annot, right=2.0cm of step1, text width=3.8cm]
    {xPert: \texttt{Perturbation[L, 2]}\\
     $\to$ \texttt{ExpandPerturbation}\\
     $\to$ truncate $O(\varepsilon^{3+})$\\
     $\to$ divide by background $\sqrt{|g_0|}$\\[3pt]
     Result: quadratic in perturbations\\
     $\to$ linear EOM after variation};

\draw[arr] (entry) -- (step1);

% ===================================================================
%  STEP 2: Matter perturbation (conditional)
% ===================================================================

\node[check, below=0.5cm of step1, text width=4.6cm] (step2)
    {\textbf{Step 2}: Matter field perturbations\\[-1pt]
     \scriptsize $A_\mu \to \bar{A}_\mu + \varepsilon\, a_\mu$, \quad
     $\phi \to \bar{\phi} + \varepsilon\, \varphi$};

\node[annot, right=2.0cm of step2, text width=3.8cm]
    {xPert's \texttt{DefTensorPerturbation}\\
     for each matter field declared\\
     in \texttt{[[linearization.}\\
     \texttt{matter\_perturbations]]} TOML.\\[3pt]
     Metric perturbation $h_{\mu\nu}$\\
     always included automatically.};

\draw[arr] (step1) -- (step2);

% ===================================================================
%  STEP 3: Fierz-Pauli fix
% ===================================================================

\node[funcbox=Purple, below=0.5cm of step2, text width=4.6cm] (step3)
    {\textbf{Step 3}: Scalar wrapper expansion\\[-1pt]
     \scriptsize Expand \texttt{Scalar[x]\^{}n} $\to$ product of renamed copies};

\node[annot, right=2.0cm of step3, text width=3.8cm]
    {Fierz--Pauli fix: prevents xPert\\
     from conflating dummy index\\
     pairs in products of Scalar[]\\
     wrappers.  Essential for correct\\
     multi-index contractions.};

\draw[arr] (step2) -- (step3);

% ===================================================================
%  STEP 4: Background substitution cascade
% ===================================================================

\node[stage=Purple, below=0.5cm of step3, text width=4.6cm] (step4)
    {\textbf{Step 4}: Background substitution};

\node[funcbox=Purple, below=0.3cm of step4, text width=3.4cm] (bg_met)
    {Metric: $g_{\mu\nu} \to \bar{g}_{\mu\nu}$};

\node[funcbox=Purple, below=of bg_met, text width=3.4cm] (bg_scalar)
    {Scalars: $\phi \to \bar{\phi}(x)$};

\node[funcbox=Purple, below=of bg_scalar, text width=3.4cm] (bg_vector)
    {Vectors: $A_\mu \to \bar{A}_\mu(x)$\\[-1pt]
     \scriptsize position-dependent profiles};

\node[annot, right=2.0cm of bg_scalar, text width=3.8cm]
    {Backgrounds must be substituted\\
     \emph{before} variation, so that\\
     VarD only varies perturbations.\\[3pt]
     Position-dependent backgrounds\\
     (e.g.\ $\bar{B}_\mu(x)$) create\\
     coordinate-dependent coefficients.};

\draw[arr] (step3) -- (step4);
\draw[arr] (step4) -- (bg_met);
\draw[arr] (bg_met) -- (bg_scalar);
\draw[arr] (bg_scalar) -- (bg_vector);

% ===================================================================
%  STEP 5: Plane-wave reduction (conditional)
% ===================================================================

\node[check, below=0.5cm of bg_vector, text width=4.6cm] (step5)
    {\textbf{Step 5}: Plane-wave reduction on $\mathcal{L}^{(2)}$\\[-1pt]
     \scriptsize zero transverse spatial derivatives\\
     \scriptsize $\to$ effective 1+1D from 3+1D};

\draw[arr] (bg_vector) -- (step5);

% ===================================================================
%  STEP 6: Torsion decomposition (conditional)
% ===================================================================

\node[check, below=0.5cm of step5, text width=4.6cm] (step6)
    {\textbf{Step 6}: Torsion curvature decomposition\\[-1pt]
     \scriptsize Riemann--Cartan $\tilde{R}_{abcd}$ $\to$\\
     \scriptsize Riemannian $R_{abcd}$ + torsion contractions};

\node[annot, right=2.0cm of step6, text width=3.8cm]
    {For Poincar\'{e} gauge theories:\\
     decompose the full curvature\\
     into GR Riemann + contorsion.\\
     Enables torsion-specific coupling\\
     analysis.};

\draw[arr] (step5) -- (step6);

% ===================================================================
%  STEP 7: Euler-Lagrange per field
% ===================================================================

\node[stage=Purple, below=0.6cm of step6, text width=4.6cm] (step7)
    {\textbf{Step 7}: Euler--Lagrange per field};

\node[stage=Pink, below=0.4cm of step7] (varloop)
    {\textbf{For each} dynamical perturbation $\delta\phi^A$};

\node[funcbox=Purple, below=of varloop] (vard)
    {\texttt{VarD[$\delta\phi^A$, CD][$\mathcal{L}^{(2)}$]}\\[-1pt]
     \scriptsize functional derivative (xAct)};

\node[funcbox=Purple, below=of vard] (canon)
    {\texttt{ToCanonical} + \texttt{ContractMetric}\\[-1pt]
     \scriptsize canonicalise abstract index form};

\draw[arr] (step6) -- (step7);
\draw[arr] (step7) -- (varloop);
\draw[arr] (varloop) -- (vard);
\draw[arr] (vard) -- (canon);

% Loop
\node[fpcheck, below=0.4cm of canon] (morefields)
    {more\\fields?};
\draw[arr] (canon) -- (morefields);
\draw[arr] (morefields.east) -- ++(1.5,0) |- (varloop.east)
    node[pos=0.25, right, font=\tiny] {yes};

% Loop background
\begin{scope}[on background layer]
    \node[loopregion, fit=(varloop)(vard)(canon)(morefields),
          label={[font=\tiny\sffamily, text=BlueLine]
          above right:per-field loop}] {};
\end{scope}

% ===================================================================
%  STEP 8: Gauge + assembly
% ===================================================================

\node[check, below=0.6cm of morefields, text width=4.6cm] (gauge)
    {\textbf{Step 8}: Gauge substitutions\\[-1pt]
     \scriptsize apply constraints (TT, Coulomb, etc.)\\
     \scriptsize to eliminate redundant components};
\draw[arr] (morefields.south) -- (gauge)
    node[midway, right, font=\tiny] {no};

\node[stage=Green, below=0.5cm of gauge] (output)
    {\textbf{Output}: \texttt{fieldEquations}\\[-1pt]
     \scriptsize list of (field, tensor EOM) pairs\\
     \scriptsize linearised, gauge-fixed};

\draw[arr] (gauge) -- (output);

% ===================================================================
%  PHASE ANNOTATIONS (left side)
% ===================================================================

\node[annot, left=3.5cm of step1, text width=2.2cm, anchor=east]
    {\textbf{Perturbation}\\theory (xPert)};
\node[annot, left=3.5cm of step4, text width=2.2cm, anchor=east]
    {\textbf{Background}\\substitution};
\node[annot, left=3.5cm of step7, text width=2.2cm, anchor=east]
    {\textbf{Variational}\\calculus (VarD)};

% Decorative brace: xPert phase
\draw[decorate, decoration={brace, amplitude=6pt, mirror},
      thick, GrayLine]
    ([xshift=-3.0cm]step1.north west) --
    ([xshift=-3.0cm]step3.south west);

% Decorative brace: background phase
\draw[decorate, decoration={brace, amplitude=6pt, mirror},
      thick, GrayLine]
    ([xshift=-3.0cm]step4.north west) --
    ([xshift=-3.0cm]step6.south west);

% Decorative brace: VarD phase
\draw[decorate, decoration={brace, amplitude=6pt, mirror},
      thick, GrayLine]
    ([xshift=-3.0cm]step7.north west) --
    ([xshift=-3.0cm]gauge.south west);

\end{tikzpicture}
\end{document}
